\section{Pole calculus for widths, line shapes, and measured response}
\label{sec:resonance-formulae}
% BEGIN CURATED SUBJECT INDEX
\index{Fredholm theory}
\index{connection coefficient}
\index{resonance!residue}
\index{Breit--Wigner approximation}
\index{Fano profile}
\index{continuum response}
% END CURATED SUBJECT INDEX
\index{resonance formula}\index{line shape}

\calculationroute
  {The pole definition of Section~\ref{sec:resonance-definition}, the
  derivative normalization of
  Section~\ref{sec:hilbert-schmidt-fredholm}, and the continuum response of
  Section~\ref{sec:continuum-response}.}
  {Derive Breit--Wigner and Fano forms as controlled approximations to a pole
  plus background, including energy-dependent self-energies, partial widths,
  thresholds, and source dependence.}
  {Given a continued connection matrix, calculate a channel residue and know
  which real-axis peak parameters are, and are not, equal to the pole data.}

A resonance formula is useful only after its numerator, denominator, sheet,
and range of approximation have been specified.  The pole position is an
analytic invariant.  A peak position, a full width at half maximum, a time
delay, and a fitted Breit--Wigner mass are real-axis summaries that coincide
with the pole parameters only under additional assumptions.  This section
derives those assumptions instead of adopting a resonance formula as a
definition.  The presentation follows the practical pole calculus emphasized
in Ref.~\cite{Kukulin:1989}, but the denominator is kept in the exact-WKB form
used throughout this book.

\subsection{Every measurement contains a source, propagation, and a detector}
\label{subsec:source-propagator-detector}
% BEGIN CURATED SUBJECT INDEX
\index{spectral measure}
\index{resolvent}
\index{resonance!residue}
% END CURATED SUBJECT INDEX

Let $\ket{F_a}$ be the state produced by an entrance channel or external
operator and let $\bra{D_b}$ represent the detected channel.  The energy-
dependent amplitude has the form
\begin{equation}
  \mathcal A_{ba}(E)
  =\bra{D_b}\Resolvent(E+\rmi0)\ket{F_a}
  \label{eq:source-detector-amplitude}
\end{equation}
up to kinematic and direct terms.  If the continued resolvent has a simple
pole $z_{\mathrm R}$, then locally
\begin{equation}
  \mathcal A_{ba}(E)
  =\vsup{\mathcal A_{ba}}{bg}(E)
  +\frac{g_b^{\mathrm R}\widetilde g_a^{\mathrm R}}
  {E-z_{\mathrm R}},
  \label{eq:amplitude-pole-background}
\end{equation}
where
\begin{align}
  g_b^{\mathrm R}
  &=\dualpair{D_b}{\Psi_{\mathrm R}},
  &\widetilde g_a^{\mathrm R}
  &=\dualpair{\widetilde\Psi_{\mathrm R}}{F_a}.
  \label{eq:pole-source-detector-couplings}
\end{align}
The residue factorizes, but its two factors depend on preparation and
detection.  A pole can be prominent in one reaction and almost invisible in
another.  Conversely, a kinematic cusp or a zero of the background amplitude
can create a peak or dip without a nearby pole.

For an inclusive response generated by an operator $O$ from a normalized
state $\ket{0}$, take $\ket{F}=O\ket{0}$ and
\begin{equation}
  S_O(E)
  =-\frac{1}{\pi}\im
  \bra{F}\Resolvent(E+\rmi0)\ket{F}.
  \label{eq:inclusive-response-pole-calculus}
\end{equation}
Positivity of the physical spectral measure makes the total $S_O(E)$
nonnegative.  A separated pole term after contour deformation need not be
positive by itself, because the continued left and right residues are not
complex conjugates and the rotated continuum completes the physical boundary
value.

\subsection{Feshbach elimination and the exact scalar denominator}
\label{subsec:feshbach-scalar-denominator}
% BEGIN CURATED SUBJECT INDEX
\index{resolvent}
\index{scattering!on-shell}
\index{resonance!residue}
\index{Feshbach projection}
\index{self-energy}
% END CURATED SUBJECT INDEX
\index{Feshbach operator}\index{self-energy}\index{wave-function renormalization}

Let $P=\ket{d}\bra{d}$ select one intrinsic doorway and $Q=\identity-P$.
Eliminating $Q$ gives
\begin{equation}
  G_{dd}(z)
  :=\bra{d}\Resolvent(z)\ket{d}
  =\frac{1}{D(z)},
  \qquad
  D(z)=z-E_d-\Sigma(z),
  \label{eq:doorway-green-denominator}
\end{equation}
with the self-energy
\begin{equation}
  \Sigma(z)
  =\bra{d}HQ(z-QHQ)^{-1}QH\ket{d}.
  \label{eq:doorway-self-energy}
\end{equation}
The resonance pole solves
\begin{equation}
  D(z_{\mathrm R})=0
  \label{eq:doorway-pole-condition}
\end{equation}
on the continued sheet.  For a simple pole,
\begin{equation}
  G_{dd}(z)
  =\frac{Z_{\mathrm R}}{z-z_{\mathrm R}}+O(1),
  \qquad
  Z_{\mathrm R}
  =\frac{1}{1-\Sigma'(z_{\mathrm R})}.
  \label{eq:doorway-wavefunction-renormalization}
\end{equation}
The quantity $Z_{\mathrm R}$ is generally complex.  It is a residue
renormalization, not a probability that must lie between zero and one.

On the upper rim of the physical cut write
\begin{equation}
  \Sigma(E+\rmi0)
  =\Delta(E)-\frac{\rmi}{2}\Gamma(E),
  \label{eq:self-energy-real-axis}
\end{equation}
where $\Delta$ is the dispersive shift.  If channel $c$ has continuum states
$\ket{E,c}$ in a fixed flux normalization, then
\begin{equation}
  \Gamma_c(E)
  =2\pi\left|\bra{E,c}H\ket{d}\right|^2,
  \qquad
  \Gamma(E)=\sum_{c\ \mathrm{open}}\Gamma_c(E)
  \label{eq:golden-rule-partial-width}
\end{equation}
in that normalization.  Equation~\eqref{eq:golden-rule-partial-width} is an
on-shell identity for the imaginary part of the physical self-energy.  At a
broad complex pole, analytic continuation and the factor
$Z_{\mathrm R}$ must be retained; simply substituting a real-energy golden-
rule width into $z_{\mathrm R}=E_d+\Delta-\rmi\Gamma/2$ is an approximation.

\subsection{The multichannel effective Hamiltonian and unitary S matrix}
\label{subsec:effective-hamiltonian-S}
% BEGIN CURATED SUBJECT INDEX
\index{cross section}
\index{scattering matrix}
\index{scattering!multichannel}
\index{resonance!residue}
\index{Breit--Wigner approximation}
\index{effective Hamiltonian}
\index{self-energy}
% END CURATED SUBJECT INDEX

Let $W_{nc}(E)$ couple intrinsic states $\ket{n}$ to open channels $c$.
With continuum states normalized per unit energy, define
\begin{align}
  \vsub{H}{eff}(E+\rmi0)
  &=H_{PP}+\Delta(E)-\rmi\pi W(E)W(E)^\dagger,
  \label{eq:weidenmuller-effective-H}\\
  S(E)
  &=\identity-2\pi\rmi\,
  W(E)^\dagger
  [E-\vsub{H}{eff}(E+\rmi0)]^{-1}
  W(E).
  \label{eq:weidenmuller-S}
\end{align}
The Hermitian principal-value term $\Delta$ and the anti-Hermitian term are
the real and imaginary parts of one self-energy.  Keeping both and using the
same $W$ makes Eq.~\eqref{eq:weidenmuller-S} unitary on the retained open
channels.  Replacing the anti-Hermitian part by an unrelated fitted width can
destroy that identity.

Suppose one isolated right--left eigenpair of the continued effective
Hamiltonian dominates.  Its contribution to
Eq.~\eqref{eq:weidenmuller-S} is
\begin{equation}
  S_{ba}(E)
  =\vsup{S_{ba}}{bg}(E)
  -\frac{\rmi\,\gamma_b^{\mathrm R}
  \widetilde\gamma_a^{\mathrm R}}
  {E-z_{\mathrm R}}
  +\cdots,
  \label{eq:multichannel-pole-S}
\end{equation}
where the channel residues contain $W$, the left and right intrinsic vectors,
and the derivative of the energy-dependent effective operator.  Only in the
narrow, slowly varying, time-reversal-invariant limit can phases be chosen so
that
\begin{equation}
  |\gamma_c^{\mathrm R}|^2\simeq\Gamma_c,
  \qquad
  \sum_c\Gamma_c\simeq\Gamma.
  \label{eq:narrow-partial-width-identification}
\end{equation}
For overlapping resonances, the residues are coherent matrices and a sum of
positive independent Breit--Wigner cross sections is not unitary.

\subsection{Derivation of the one-channel Breit--Wigner factor}
\label{subsec:breit-wigner-derivation}
% BEGIN CURATED SUBJECT INDEX
\index{phase shift}
\index{Wigner time delay}
\index{scattering matrix}
\index{resonance!residue}
\index{Breit--Wigner approximation}
\index{self-energy}
% END CURATED SUBJECT INDEX
\index{Breit--Wigner formula}\index{Wigner time delay}

For one elastic channel, real analyticity and unitarity imply that an isolated
simple pole $z_{\mathrm R}=E_*-\rmi\Gamma/2$ is accompanied by a zero at
$z_{\mathrm R}^*$.  Factor the scattering matrix as
\begin{equation}
  S(E)
  =\rme^{2\rmi\vsub{\delta}{bg}(E)}
  \frac{E-z_{\mathrm R}^*}{E-z_{\mathrm R}}
  \vsub{S}{far}(E).
  \label{eq:unitary-pole-factor}
\end{equation}
If $\vsub{\delta}{bg}$ and $\vsub{S}{far}$ vary negligibly over an
interval of several widths, absorb them into a constant phase.  The resonant
partial-wave amplitude then has denominator
\begin{equation}
  E-E_*+\frac{\rmi}{2}\Gamma,
  \label{eq:breit-wigner-amplitude-denominator}
\end{equation}
and its modulus squared contains
\begin{equation}
  \frac{1}{(E-E_*)^2+\Gamma^2/4}.
  \label{eq:breit-wigner-lorentzian}
\end{equation}
The pole definition came first; the Lorentzian followed after freezing the
background, residue, phase space, and self-energy.

For the resonant phase shift one may take
\begin{equation}
  \frac{E-z_{\mathrm R}^*}{E-z_{\mathrm R}}
  =\rme^{2\rmi\delta_{\mathrm R}(E)},
  \qquad
  \frac{\rmd\delta_{\mathrm R}}{\rmd E}
  =\frac{\Gamma/2}
  {(E-E_*)^2+\Gamma^2/4}.
  \label{eq:resonant-phase-derivative}
\end{equation}
Thus the Wigner delay contributed by the isolated factor is
\begin{equation}
  \tau_{\mathrm R}(E)
  =\frac{\hbar\Gamma}
  {(E-E_*)^2+\Gamma^2/4},
  \qquad
  \tau_{\mathrm R}(E_*)=\frac{4\hbar}{\Gamma}.
  \label{eq:breit-wigner-time-delay}
\end{equation}
The peak delay differs by a numerical factor from the exponential lifetime
$\hbar/\Gamma$ because it measures the energy derivative of a round-trip
scattering phase, not the survival time of a normalized unstable packet.

\subsection{Background interference and the Fano profile}
\label{subsec:fano-from-pole}
% BEGIN CURATED SUBJECT INDEX
\index{resonance!residue}
\index{Fano profile}
% END CURATED SUBJECT INDEX
\index{Fano profile}\index{background interference}

Keep the coherent background in
Eq.~\eqref{eq:amplitude-pole-background}.  Over a narrow interval approximate
$\vsup{\mathcal A}{bg}=A_0$ and the residue by $r_0$.  With
\begin{equation}
  \epsilon=\frac{2(E-E_*)}{\Gamma},
  \label{eq:fano-reduced-energy}
\end{equation}
the amplitude can be written, after extracting a constant phase and scale,
as
\begin{equation}
  \mathcal A(E)
  =A_0\frac{\epsilon+q}{\epsilon+\rmi}.
  \label{eq:fano-amplitude}
\end{equation}
When time-reversal and channel conventions make $q$ real, the intensity is
\begin{equation}
  |\mathcal A(E)|^2
  =|A_0|^2\frac{(\epsilon+q)^2}{1+\epsilon^2}.
  \label{eq:fano-profile}
\end{equation}
A finite $q$ shifts the maximum away from $E_*$ and creates a zero at
$\epsilon=-q$.  With absorption or several coherent backgrounds, $q$ is
complex and the exact zero is filled.  The asymmetric line shape is still a
pole--background interference effect; its maximum is not a new definition of
the resonance energy.

\subsection{Threshold laws and why broad resonances are not Lorentzians}
\label{subsec:threshold-resonance-formula}
% BEGIN CURATED SUBJECT INDEX
\index{threshold law}
\index{scattering!Coulomb}
\index{nuclear breakup}
% END CURATED SUBJECT INDEX
\index{threshold law}\index{Wigner threshold law}

Near a neutral two-body threshold, phase space and the centrifugal barrier
give the Wigner law
\begin{equation}
  \Gamma_l(E)
  \propto k^{2l+1},
  \qquad
  k=\sqrt{2\mu(E-\vsub{E}{th})}/\hbar.
  \label{eq:wigner-width-law}
\end{equation}
The square root is the same branch that generates the energy sheets.  A more
faithful one-level denominator is therefore
\begin{equation}
  D(E)=E-E_d-\Delta(E)+\frac{\rmi}{2}\Gamma(E),
  \label{eq:energy-dependent-resonance-denominator}
\end{equation}
not a constant-width expression.  The dispersive part $\Delta(E)$ has the
corresponding nonanalyticity.  If the pole is close to threshold or
$\Gamma$ is comparable with its distance from other singularities, the
maximum of $|D(E)|^{-2}$ need not occur at $\re z_{\mathrm R}$, and its
half-maximum width need not equal $-2\im z_{\mathrm R}$.

For a charged channel the penetrability contains Coulomb functions and an
essential threshold factor related to $\rme^{-2\pi\eta}$.  The ordinary
power law must be replaced by the Coulomb-modified one.  For a genuine
three-body breakup the threshold density and final-state interactions are
multidimensional; reducing them to a constant ``three-body width'' can erase
the very correlations measured in an invariant-mass spectrum.

\subsection{Connection-matrix residue: the exact-WKB resonance formula}
\label{subsec:connection-matrix-resonance-formula}
% BEGIN CURATED SUBJECT INDEX
\index{Fredholm theory}
\index{connection coefficient}
\index{boundary condition!incoming and outgoing}
\index{resonance!residue}
\index{Breit--Wigner approximation}
\index{exact WKB}
\index{exceptional point}
\index{exceptional point!defectiveness}
% END CURATED SUBJECT INDEX
\index{connection matrix}\index{exact WKB!resonance residue}

In one channel the exact-WKB scattering amplitude can be arranged as
\begin{equation}
  S(E)=\frac{\mathcal N(E)}{\Cincoming(E)}.
  \label{eq:wkb-S-numerator-denominator}
\end{equation}
If $\Cincoming(z_{\mathrm R})=0$ is simple, then
\begin{equation}
  \Res_{E=z_{\mathrm R}}S(E)
  =\frac{\mathcal N(z_{\mathrm R})}
  {\partial_E\Cincoming(z_{\mathrm R})}.
  \label{eq:wkb-S-residue-formula}
\end{equation}
The numerator is fixed by the remaining connection data and flux convention.
The denominator is a global derivative: both the perturbative WKB periods and
the exponentially small Stokes multipliers must be differentiated.

For $N$ channels, let $\Cincoming(E)$ map the independent
interior or left data to prohibited incoming amplitudes.  At a simple
resonance choose
\begin{equation}
  \Cincoming(z_{\mathrm R})u=0,
  \qquad
  v^\dagger\Cincoming(z_{\mathrm R})=0.
  \label{eq:connection-left-right-null}
\end{equation}
Then
\begin{equation}
  \Cincoming(E)^{-1}
  =\frac{u v^\dagger}
  {(E-z_{\mathrm R})
   v^\dagger\Cincoming'(z_{\mathrm R})u}
  +\order(1).
  \label{eq:connection-inverse-pole}
\end{equation}
Multiplying Eq.~\eqref{eq:connection-inverse-pole} by the matrices that
couple the source to $u$ and $v$ to the detector gives every channel residue.
This is the same finite-dimensional obstruction obtained from analytic
Fredholm theory, now calculated by exact-WKB transport.

At an exceptional point the scalar denominator in
Eq.~\eqref{eq:connection-inverse-pole} vanishes together with the first
derivative or the null space becomes defective.  The Laurent expansion then
contains a higher-order pole and associated vectors.  Applying a simple
Breit--Wigner formula at that point discards the defining Jordan data; Stage
VII derives the correct replacement.

\subsection{Complex scaling separates a pole term without discarding the cut}
\label{subsec:resonance-formula-csm}
% BEGIN CURATED SUBJECT INDEX
\index{spectrum!point}
\index{resonance!residue}
\index{complex scaling}
\index{pseudostate}
% END CURATED SUBJECT INDEX

Using the extended complex-scaled resolution, the tested response becomes
\begin{equation}
  S_O(E)
  \approx-\frac{1}{\pi}\im
  \sum_\nu
  \frac{
  \bra{F}U^{-1}(\theta)\ket{\Phi_\nu^\theta}
  \bra{\widetilde\Phi_\nu^\theta}U(\theta)\ket{F}}
  {E-E_\nu^\theta}.
  \label{eq:resonance-formula-csm-response}
\end{equation}
Stable isolated eigenvalues represent exposed poles; the eigenvalues aligned
with rotated rays discretize the deformed cuts.  Selecting one stable term is
a useful diagnostic of its source-dependent residue.  The remaining terms
are not numerical contamination: they reconstruct nonresonant background,
threshold cusps, and the interference needed for the physical line shape.

This point is essential for the broad $2_2^+$ state of $^6$He in Laboratory
VI-A.  The pole may be isolated in the complex-scaled spectrum while the
observed energy distribution still depends strongly on nonresonant continuum
terms.  ``The resonance contribution'' must therefore name the Riesz pole
term and the chosen source, not an arbitrary window of pseudostate energies.

\subsection{A line-shape calculation that can be reproduced}
\label{subsec:line-shape-protocol}
% BEGIN CURATED SUBJECT INDEX
\index{scattering!multichannel}
\index{connection coefficient}
\index{boundary condition!incoming and outgoing}
\index{resonance!residue}
\index{Breit--Wigner approximation}
\index{Fano profile}
% END CURATED SUBJECT INDEX

For any proposed resonance formula, record the following in this order:
\begin{enumerate}[label=\arabic*.]
  \item the continued variable and sheet on which $z_{\mathrm R}$ is found;
  \item the analytic denominator and the test for a simple or multiple zero;
  \item the left--right derivative normalization and channel residue;
  \item the source and detector operators used to construct the numerator;
  \item every open-channel threshold and the energy dependence of its
  penetrability;
  \item the background amplitude retained coherently with the pole;
  \item the energy interval over which that background and the residue were
  approximated as constant;
  \item a comparison of the pole parameters with the peak, half-maximum, and
  time-delay parameters rather than an identification by notation.
\end{enumerate}

\begin{checkpointbox}
The reader should now be able to start from $D(z)=z-E_d-\Sigma(z)$, derive
$Z_{\mathrm R}=[1-\Sigma'(z_{\mathrm R})]^{-1}$, obtain the Breit--Wigner
and Fano profiles only after declaring their approximations, and calculate a
multichannel residue from
$v^\dagger\Cincoming'(z_{\mathrm R})u$.
\end{checkpointbox}
